\section{Related Work}
\label{sec:related}

\paragraph{Code-generation benchmarks.}
HumanEval~\citep{chen2021codex}, MBPP~\citep{austin2021mbpp}, and
APPS~\citep{hendrycks2021apps} evaluate algorithmic Python on isolated
functions; BigCodeBench~\citep{zhuo2024bigcodebench} extends the paradigm to
diverse function calls and complex instructions.
SWE-Bench~\citep{jimenez2024swebench} and its agent-driven successor
SWE-Agent~\citep{yang2024swe} target repository-level bug fixing within
open-source Python projects; RepoBench~\citep{liu2024repobench} benchmarks
cross-file code completion; LiveCodeBench~\citep{livecodebench2024} evaluates
on continuously updated competitive-programming problems to reduce data
contamination.
These benchmarks share our emphasis on deterministic scoring---unit tests,
exact match, AST equivalence---but address general-purpose programming.
None cover SAS or R in regulated settings, domain-specific data standards
such as CDISC SDTM/ADaM~\citep{cdisc_sdtm_ig,cdisc_adam_ig}, or the
documentation and validation artefacts required in regulatory submissions.
ClinProg-Bench fills this gap by grounding all tasks in a real
clinical-trial substrate and scoring against CDISC-conforming gold outputs.

\paragraph{Domain-specific code benchmarks.}
DS-1000~\citep{lai2023ds1000} evaluates data-science code generation across
NumPy, Pandas, and TensorFlow with a ``strict'' mode that rejects syntactic
variants; Spider~\citep{yu2018spider} benchmarks text-to-SQL semantic parsing
against gold query sets; CrossCodeEval~\citep{ding2023crosscodeeval} measures
cross-file code completion across six programming languages;
CRUXEval~\citep{gu2024cruxeval} probes code reasoning and execution
understanding.
These works share our emphasis on grounded, deterministic evaluation but
address different domains without regulatory constraints.
A key distinction is that ClinProg-Bench tasks produce outputs that must
conform to externally imposed standards (CDISC controlled terminology, ADaM
variable-naming conventions) rather than satisfying only functional
correctness.
A program that produces numerically correct output may still fail evaluation
if variable names, metadata, or documentation do not comply with the
submission standard---a qualitatively different acceptance criterion.

\paragraph{Biomedical and clinical evaluations.}
MedQA~\citep{jin2021medqa}, PubMedQA~\citep{jin2019pubmedqa}, and
MedAlign~\citep{fleming2024medalign} probe clinical and biomedical reasoning
over natural language and structured text.
These benchmarks evaluate text understanding and clinical decision-making but
do not assess code generation, data-standards conformance, or the software
engineering pipeline that produces regulatory datasets and tables.
Industry standards bodies---CDISC~\citep{cdisc_sdtm_ig,cdisc_adam_ig}, PHUSE,
and the ICH~\citep{ich_e6}---define the data models, controlled terminologies,
and process frameworks (GCP) that clinical programmers work within, but they
provide substrate and guidance rather than benchmarking infrastructure.
ClinProg-Bench bridges this gap: it is the first public benchmark that
evaluates AI systems on the full clinical-programming lifecycle (code
generation, review, specification extraction, documentation, debugging)
against a CDISC-conforming substrate released under an open license.

\paragraph{Agentic benchmarks.}
GAIA~\citep{mialon2023gaia} evaluates general-purpose AI assistants on tasks
requiring web browsing, tool use, and multi-step reasoning;
OSWorld~\citep{xie2024osworld} benchmarks multimodal agents in real desktop
environments; WebArena~\citep{webarena2023} tests end-to-end web interaction;
AgentBench~\citep{agentbench2023} provides a multi-environment evaluation
suite across operating systems, databases, and web tasks.
These benchmarks demonstrate the value of stateful, multi-step evaluation
protocols for agentic systems.
ClinProg-Bench adapts this philosophy to the regulated software domain: its
stateful runner contract (\S\ref{sec:benchmark}) supports multi-agent systems
that can read specifications, write programs, inspect logs, and
iterate---mirroring the real clinical-programming workflow.
Unlike general agentic benchmarks, ClinProg-Bench evaluates domain-specific
competence with deterministic oracles, avoiding the reliability and
reproducibility concerns of LLM-as-a-judge scoring.
